% Standalone problem document, generated from the adjacent README.md.
% Compile with XeLaTeX. Mathematical content is maintained in Markdown.
\documentclass[11pt,a4paper]{article}
\usepackage[fontsize=10.5pt]{scrextend}
\usepackage[margin=22mm,headheight=15pt,headsep=9mm,footskip=11mm]{geometry}
\usepackage{fontspec}
\setmainfont{texgyrepagella}[Extension=.otf,UprightFont=*-regular,BoldFont=*-bold,ItalicFont=*-italic,BoldItalicFont=*-bolditalic]
\setsansfont{texgyreheros}[Extension=.otf,UprightFont=*-regular,BoldFont=*-bold,ItalicFont=*-italic,BoldItalicFont=*-bolditalic]
\setmonofont{texgyrecursor}[Extension=.otf,UprightFont=*-regular,BoldFont=*-bold,ItalicFont=*-italic,BoldItalicFont=*-bolditalic,Scale=MatchLowercase]
\usepackage{amsmath,amssymb}
\usepackage{unicode-math}
\setmathfont{texgyrepagella-math.otf}
\usepackage{xcolor,microtype,enumitem,needspace,fancyhdr}
\usepackage{bookmark}
\definecolor{ink}{HTML}{17344B}
\definecolor{accent}{HTML}{197D80}
\definecolor{muted}{HTML}{596773}
\hypersetup{colorlinks=true,linkcolor=accent,urlcolor=accent,pdftitle={NR-04 - The
nonnegative rank of the nine-point distance
matrix},pdfauthor={Open Problems in Numerical Linear Algebra}}
\urlstyle{same}
\setlength{\parindent}{0pt}
\setlength{\parskip}{5pt plus 1pt minus 1pt}
\linespread{1.04}
\setlength{\emergencystretch}{3em}
\widowpenalty=10000
\clubpenalty=10000
\displaywidowpenalty=10000
\allowdisplaybreaks[1]
\setlist{nosep,leftmargin=1.5em}
\providecommand{\tightlist}{\setlength{\itemsep}{0pt}\setlength{\parskip}{0pt}}
\providecommand{\pandocbounded}[1]{#1}
\pagestyle{fancy}
\fancyhf{}
\fancyhead[L]{\sffamily\footnotesize\color{muted}OPEN PROBLEMS IN NUMERICAL LINEAR ALGEBRA}
\fancyhead[R]{\sffamily\bfseries\footnotesize\color{accent}NR-04}
\fancyfoot[L]{\parbox[b]{0.9\textwidth}{\sffamily\fontsize{8}{10}\selectfont\color{muted}Editorial ratings; a bounded search cannot certify that a problem remains unsolved.\\Literature check: 2026-09-11}}
\fancyfoot[R]{\sffamily\footnotesize\color{muted}\thepage}
\renewcommand{\headrulewidth}{0.3pt}
\renewcommand{\headrule}{\hbox to\headwidth{\color{accent}\leaders\hrule height \headrulewidth\hfill}}
\makeatletter
\renewcommand{\section}{\Needspace{5\baselineskip}\@startsection{section}{1}{0pt}{12pt plus 2pt minus 2pt}{4pt}{\sffamily\bfseries\large\color{ink}}}
\renewcommand{\subsection}{\Needspace{4\baselineskip}\@startsection{subsection}{2}{0pt}{9pt plus 1pt minus 1pt}{3pt}{\sffamily\bfseries\normalsize\color{ink}}}
\makeatother
\setcounter{secnumdepth}{0}
\begin{document}
{\sffamily\small\color{muted}Nonnegative and positive
factorizations\par}
\vspace{3pt}
{\raggedright\hyphenpenalty=10000\exhyphenpenalty=10000\sffamily\bfseries\fontsize{21}{24}\selectfont\color{ink}The
nonnegative rank of the nine-point distance matrix\par}
\vspace{7pt}
{\sffamily\small\textbf{Difficulty:} hard\quad\textbf{Importance:} interesting
to specialist\par}
{\sffamily\small\bfseries\color{accent}Status: Solved\par}
\vspace{4pt}
{\color{accent}\hrule height 0.8pt}
\vspace{6pt}
\textbf{Rating rationale:} Hard because a single small exact
factorization must be constructed or excluded using sharp
nonnegative-rank tools; specialist importance reflects its role as a
named distance-matrix benchmark.

\textbf{Area:} exact factorization of Euclidean distance matrices

\subsection{Resolution --- 2026-09-11}\label{resolution-2026-09-11}

\textbf{Author:} Matthew J. Colbrook, Department of Applied Mathematics
and Theoretical Physics, University of Cambridge. \textbf{Independent
Codex-agent review: PASS for the exact target.}

The nine-point matrix \(D_{ij}=(i-j)^2\) has nonnegative rank seven, so
no exact six-term nonnegative factorization exists. A polygon-contact
argument applied to both factors, together with Sylvester's rank
inequality, proves the lower bound; an explicit seven-term integer
factorization proves the upper bound.

The complete target is resolved. Its former difficulty rating is
historical; the original statement, references and dated audits remain
below.

\textbf{Primary reference:}
\href{https://github.com/ajt60gaibb/OpenProblemsInNLA/blob/main/references/colbrook-factorization-2026-09-11/manuscripts/NR-04_nine_point_distance.pdf}{complete
authored PDF},
\href{https://github.com/ajt60gaibb/OpenProblemsInNLA/blob/main/references/colbrook-factorization-2026-09-11/manuscripts/NR-04_nine_point_distance.tex}{standalone
TeX}, \textbf{Theorem 1, with Theorem 4 for the lower bound}.
\href{https://github.com/ajt60gaibb/OpenProblemsInNLA/blob/main/references/colbrook-factorization-2026-09-11/verification/reviews/NR-04-review.md}{Independent
proof review} ·
\href{https://github.com/ajt60gaibb/OpenProblemsInNLA/blob/main/references/colbrook-factorization-2026-09-11/README.md}{Authorship
and submission record}. Verification is independent agent review, not
external human peer review or formal certification.

\subsection{Context and notation}\label{context-and-notation}

All factorizations are over the real numbers. For
\(X\in\mathbb R_{\ge0}^{m\times n}\), define

\[
\operatorname{rank}_+(X)=\min\{r\ge0:X=WH,\quad
W\in\mathbb R_{\ge0}^{m\times r},\ H\in\mathbb R_{\ge0}^{r\times n}\}.
\]

\subsection{Problem statement}\label{problem-statement}

Let \(D\in\mathbb R_{\ge0}^{9\times9}\) have entries \(D_{ij}=(i-j)^2\)
for \(1\le i,j\le9\).

\subsubsection{Question}\label{question}

Does there exist an exact factorization \(D=WH\) with
\(W\in\mathbb R_{\ge0}^{9\times6}\) and
\(H\in\mathbb R_{\ge0}^{6\times9}\)? Equivalently, is
\(\operatorname{rank}_+(D)=6\) or \(7\)?

\subsection{References}\label{references}

Baeckelant, Vandaele, and Gillis,
\href{https://arxiv.org/html/2605.14058v2}{\emph{Computing Lower Bounds
on the Nonnegative Rank via Non-Convex Optimization Solvers}}, Appendix
A.1, Table 6 and its final paragraph, expressly isolates this case.
Pavel Hrubeš, \emph{On the nonnegative rank of distance matrices},
Information Processing Letters \textbf{112} (2012), 457--461, supplies
the upper-bound construction cited there.

\subsection{Status check --- 2026-09-10}\label{status-check-2026-09-10}

Rechecked \href{https://arxiv.org/html/2605.14058v2}{Baeckelant et
al.~v2, Appendix A.1, Table 6 and final paragraph}, and searched for
later nine-point distance-matrix factorizations. The July 2026 text
explicitly retains the gap between six and seven. Its new exact results
for orders eleven and twelve do not settle order nine. No six-factor
construction or seven-factor lower bound was located.
\end{document}
